\documentclass[12pt]{ximera}
\title{Homework 2: Modified Borda and Plurality-with-Elimination}
\begin{document}
\begin{abstract}
Sections 1.3--1.4: a modified Borda count with custom point values, Plurality-with-Elimination
round by round, and what happens to the same election when ballots are truncated to
the voters' top three choices.
\end{abstract}
\maketitle

\section*{Sections 1.3--1.4: Modified Borda and Plurality-with-Elimination}

\begin{problem}
Table 1 shows the preference schedule for Best Album Covers at the Lakefront Music
Fest. The candidates are Al, Bruce, Cheryl, and Doreen (A, B, C, D).

\begin{center}
\begin{tabular}{c|cccccc}
Preference & 12 & 11 & 8 & 5 & 4 & 2 \\\hline
1st & A & C & D & B & B & A \\
2nd & B & B & A & C & D & D \\
3rd & C & D & C & D & A & C \\
4th & D & A & B & A & C & B
\end{tabular}

\smallskip
Table 1. Preference schedule for Best Album Covers.
\end{center}

The festival uses a modified Borda count where 1st place votes are worth $7$ points,
2nd place $5$ points, 3rd place $2$ points, and 4th place $1$ point.

\begin{prompt}
Find the complete ranking of the four best album covers, with each score.

1st: $\answer{B}$ with $\answer{188}$ \qquad
2nd: $\answer{A}$ with $\answer{162}$

\smallskip
3rd: $\answer{C}$ with $\answer{150}$ \qquad
4th: $\answer{D}$ with $\answer{130}$

\begin{hint}
The method is the same as an ordinary Borda count --- only the point values change.
Go candidate by candidate through all six columns.
\end{hint}

\begin{solution}
Taking B as a worked example:
\[
\begin{array}{lll}
\text{12 voters: B is 2nd} & 5\text{ pts} & 60\\
\text{11 voters: B is 2nd} & 5\text{ pts} & 55\\
\text{8 voters: B is 4th}  & 1\text{ pt}  & 8\\
\text{5 voters: B is 1st}  & 7\text{ pts} & 35\\
\text{4 voters: B is 1st}  & 7\text{ pts} & 28\\
\text{2 voters: B is 4th}  & 1\text{ pt}  & 2
\end{array}
\]
for a total of $188$. Doing the same for the others gives
\[ \text{B } 188 > \text{A } 162 > \text{C } 150 > \text{D } 130. \]
As a check, the four scores total $630 = 42\cdot 15$, since each of the $42$ ballots
hands out $7+5+2+1=15$ points.

Note that B wins despite having only $9$ first-place votes to A's $14$ --- the modified
weights still reward B's many second-place finishes.
\end{solution}
\end{prompt}
\end{problem}

\begin{problem}
Table 2 gives a preference schedule for election to the British Cheese Board.

\begin{center}
\begin{tabular}{c|cccccc}
Preference & 22 & 15 & 11 & 10 & 6 & 2 \\\hline
1st & A & E & C & E & B & A \\
2nd & B & B & B & D & C & E \\
3rd & C & C & D & C & A & D \\
4th & D & D & E & A & D & B \\
5th & E & A & A & B & E & C
\end{tabular}

\smallskip
Table 2. Preference schedule for the British Cheese Board.
\end{center}

\begin{prompt}
Using the Plurality-with-Elimination method, find the winner and the complete ranking.
Record the vote totals in each round.

\textbf{Round 1:}
A: $\answer{24}$ \quad B: $\answer{6}$ \quad C: $\answer{11}$ \quad D: $\answer{0}$ \quad E: $\answer{25}$

\smallskip
Eliminated in round 1: $\answer{D}$

\smallskip
\textbf{Round 2:} A: $\answer{24}$ \quad B: $\answer{6}$ \quad C: $\answer{11}$ \quad E: $\answer{25}$ \qquad
Eliminated: $\answer{B}$

\smallskip
\textbf{Round 3:} A: $\answer{24}$ \quad C: $\answer{17}$ \quad E: $\answer{25}$ \qquad
Eliminated: $\answer{C}$

\smallskip
\textbf{Round 4:} A: $\answer{30}$ \quad E: $\answer{36}$

\smallskip
Winner: $\answer{E}$

\begin{hint}
When a candidate is eliminated, their votes transfer to whichever remaining candidate
is highest on each of those ballots. Nobody's ballot is discarded here, so every round
should still total $66$.
\end{hint}

\begin{solution}
There are $66$ voters.

\textbf{Round 1.} First-place votes: A $= 22+2 = 24$, E $= 15+10 = 25$, C $= 11$,
B $= 6$, D $= 0$. D has the fewest and is eliminated --- D receives no first-place
votes at all.

\textbf{Round 2.} Removing D changes nobody's top choice, so the counts are unchanged:
A $24$, E $25$, C $11$, B $6$. B is eliminated.

\textbf{Round 3.} The $6$ voters whose first choice was B had C second, so those $6$
votes go to C: C $= 11+6 = 17$. Now A $24$, E $25$, C $17$. C is eliminated.

\textbf{Round 4.} C's $17$ votes split: the $11$-voter column ranked D, E, A after C,
and with D gone those go to E; the $6$-voter column ranked A next among survivors, so
those go to A. That gives
\[ \text{A} = 24+6 = 30, \qquad \text{E} = 25+11 = 36. \]
E wins.

The complete ranking, in reverse order of elimination, is
\[ \text{E},\ \text{A},\ \text{C},\ \text{B},\ \text{D}. \]
\end{solution}
\end{prompt}
\end{problem}

\begin{problem}
Suppose the voters in the British Cheese Board election were only asked about their top
three choices, giving the schedule in Table 3 (with no other change).

\begin{center}
\begin{tabular}{c|cccccc}
Preference & 22 & 15 & 11 & 10 & 6 & 2 \\\hline
1st & A & E & C & E & B & A \\
2nd & B & B & B & D & C & E \\
3rd & C & C & D & C & A & D
\end{tabular}

\smallskip
Table 3. Truncated preference schedule.
\end{center}

\begin{prompt}
Using truncated Plurality-with-Elimination, find the winner and complete ranking,
noting any exhausted ballots.

\textbf{Round 1:} A: $\answer{24}$ \quad B: $\answer{6}$ \quad C: $\answer{11}$ \quad
D: $\answer{0}$ \quad E: $\answer{25}$ \qquad Eliminated: $\answer{D}$

\smallskip
\textbf{Round 2:} A: $\answer{24}$ \quad B: $\answer{6}$ \quad C: $\answer{11}$ \quad
E: $\answer{25}$ \qquad Eliminated: $\answer{B}$

\smallskip
\textbf{Round 3:} A: $\answer{24}$ \quad C: $\answer{17}$ \quad E: $\answer{25}$
\qquad Eliminated: $\answer{C}$

\smallskip
\textbf{Round 4:} A: $\answer{30}$ \quad E: $\answer{25}$ \quad
exhausted ballots: $\answer{11}$

\smallskip
Winner: $\answer{A}$

\begin{hint}
A ballot becomes \emph{exhausted} when every candidate it listed has been eliminated ---
it then stops counting entirely, so later rounds may total fewer than $66$ votes.
\end{hint}

\begin{solution}
The first three rounds are identical to the full election, because the eliminations of
D, B, and C all resolve within the voters' top three choices.

\textbf{Round 4} is where truncation bites. With B, C, and D all eliminated, look at
what each column has left:
\begin{itemize}
\item $22$ voters (A, B, C): A survives $\to$ A.
\item $15$ voters (E, B, C): E survives $\to$ E.
\item $11$ voters (C, B, D): \emph{all three} are eliminated $\to$ ballot exhausted.
\item $10$ voters (E, D, C): E survives $\to$ E.
\item $6$ voters (B, C, A): A survives $\to$ A.
\item $2$ voters (A, E, D): A survives $\to$ A.
\end{itemize}
So A $= 22+6+2 = 30$, E $= 15+10 = 25$, and $11$ ballots are exhausted. Only $55$
votes still count, and \textbf{A wins}.

The complete ranking is A, E, C, B, D.

This is the whole point of the problem: truncating the ballots \emph{changed the
winner} from E to A. The $11$ voters who ranked C first had E fourth on their full
ballots, and in the complete election those votes carried E to victory. Once those
voters were only allowed three choices, their ballots fell out of the count entirely
and A won instead.
\end{solution}
\end{prompt}
\end{problem}

\end{document}
